% Owner: Anshika (Member 2)
%=====================================================================
\section{Proposed Solution}
%=====================================================================
\subsection{Phase I: Harvard system and shared-memory baseline}
Fig.~\ref{fig:block1} shows the two Phase~I systems. \sig{harvard\_top} connects a single-cycle
processor (\sig{cpu\_core}) to a 256-word instruction ROM and, through an address decoder and a
read/write controller, to a 256-word data RAM and a memory-mapped output register at
\texttt{0x400}. Any other data address is invalid: a store is ignored, a load returns 0 and
\sig{d\_fault} is raised. The two ports share no bus and no storage, so a fetch and a load or
store complete in the same cycle. The baseline \sig{vn\_top} reuses the same register file, ALU,
immediate generator and decoder in \sig{cpu\_core\_vn}, but reaches one 512-word memory (code in
words 0--255, data in words 256--511) through a two-state controller (Fig.~\ref{fig:fsm}).
Loads and stores therefore take a \textsc{FetchExec} cycle and a \textsc{Mem} cycle.

\begin{figure*}[t]
  \centering
  \resizebox{0.80\textwidth}{!}{% Block diagrams of the two systems (Fig. 1)
\begin{tikzpicture}[
  font=\scriptsize,
  blk/.style={draw, rounded corners=2pt, align=center, minimum height=0.8cm, fill=blue!7},
  cpu/.style={draw, rounded corners=2pt, align=center, fill=orange!12, minimum width=1.9cm},
  ctr/.style={draw, rounded corners=2pt, align=center, fill=gray!12, minimum height=0.5cm},
  arr/.style={-{Latex[length=1.6mm]}, thick},
  lbl/.style={font=\tiny\ttfamily, fill=white, inner sep=1pt}]

% ---------------- (a) Harvard ----------------
\node[font=\small\bfseries] at (2.6,3.0) {(a) Harvard: \texttt{harvard\_top}};
\node[cpu, minimum height=3.4cm] (core) at (0,0.4) {\texttt{cpu\_core}\\[2pt]single-cycle\\RV32I subset};
\node[blk, minimum width=2.4cm] (imem) at (4.2,1.6) {\texttt{instr\_mem}\\256$\times$32 ROM};
\node[blk, minimum width=2.4cm] (dec)  at (4.2,0.0) {\texttt{addr\_decoder}\\+ \texttt{rw\_ctrl}};
\node[blk, minimum width=2.4cm] (dmem) at (4.2,-1.4) {\texttt{data\_mem}\\256$\times$32 RAM};
\node[ctr] (out) at (6.3,0.0) {\texttt{out\_port}};
\draw[arr] ([yshift=4pt]core.east |- imem.west) -- node[lbl,above]{i\_addr} ([yshift=4pt]imem.west);
\draw[arr] ([yshift=-4pt]imem.west) -- node[lbl,below]{i\_rdata} ([yshift=-4pt]core.east |- imem.west);
\draw[arr] ([yshift=4pt]core.east |- dec.west) -- node[lbl,above]{d\_addr, d\_wdata} ([yshift=4pt]dec.west);
\draw[arr] ([yshift=-4pt]dec.west) -- node[lbl,below]{d\_rdata} ([yshift=-4pt]core.east |- dec.west);
\draw[arr] (dec.south) -- node[lbl,right]{ram\_we} (dmem.north);
\draw[arr] (dec.east) -- (out.west);
\draw[dashed, gray] (2.7,0.8) -- (5.6,0.8) node[right, font=\tiny, text=gray]{separate ports};
\node[ctr, minimum width=3.6cm] (pc) at (2.4,-2.5) {\texttt{perf\_counters}};
\draw[arr] (core.south) |- (pc.west);

% ---------------- (b) Shared memory ----------------
\begin{scope}[xshift=9.6cm]
\node[font=\small\bfseries] at (2.0,3.0) {(b) Baseline: \texttt{vn\_top}};
\node[cpu, minimum height=3.4cm] (vcore) at (0,0.4) {\texttt{cpu\_core\_vn}\\[2pt]2-state FSM\\FETCH\_EXEC\\MEM};
\node[blk, minimum width=2.6cm] (vdec) at (4.0,1.4) {address mux\\(+0x400 for data)};
\node[blk, minimum width=2.6cm] (smem) at (4.0,-0.8) {\texttt{shared\_mem}\\512$\times$32\\code 0--255, data 256--511};
\draw[arr] (vcore.east |- vdec.west) -- node[lbl,above]{m\_addr} (vdec.west);
\draw[arr] (vdec.south) -- node[lbl,right]{single port} (smem.north);
\draw[arr] (smem.west) -- node[lbl,below]{m\_rdata} (vcore.east |- smem.west);
\node[ctr, minimum width=3.6cm] (vpc) at (2.4,-2.5) {\texttt{perf\_counters}};
\draw[arr] (vcore.south) |- (vpc.west);
\end{scope}
\end{tikzpicture}
}
  \caption{Phase I. (a) Harvard system: the processor reaches the instruction ROM and the data
  side through separate ports. (b) Shared-memory baseline: fetches and data accesses use one
  port to a single memory.}
  \label{fig:block1}
\end{figure*}

\begin{figure}[t]
  \centering
  \resizebox{0.62\columnwidth}{!}{% Two-state controller of the shared-memory baseline (Fig. 3)
\begin{tikzpicture}[
  font=\scriptsize,
  st/.style={draw, circle, thick, fill=orange!12, minimum size=1.55cm, align=center},
  w/.style={-{Latex[length=1.6mm]}, thick}]
\node[st] (fe) at (0,0) {FETCH\_\\EXEC\\[1pt]{\tiny bus = PC}};
\node[st] (me) at (3.6,0) {MEM\\[1pt]{\tiny bus = data addr}};
\draw[w] (-1.6,0) node[left]{rst} -- (fe);
\draw[w] (fe) to[loop above, looseness=6] node[above, align=center]{ADD/SUB/AND/OR/ADDI,\\BEQ/BNE, EBREAK: 1 cycle} (fe);
\draw[w] (fe) to[bend left=28] node[above, align=center]{LW / SW:\\latch instruction} (me);
\draw[w] (me) to[bend left=28] node[below, align=center]{access done,\\PC $\leftarrow$ PC+4} (fe);
\end{tikzpicture}
}
  \caption{Two-state controller of the single-bus core (Phase I, kept in Phase II).}
  \label{fig:fsm}
\end{figure}

\subsection{Phase II architecture}
Phase~II keeps both processors and puts a memory hierarchy behind them
(Fig.~\ref{fig:block2}). Every Phase~II system uses the same lower-level memory,
\sig{main\_mem}, a single-port memory that holds code at physical addresses
\texttt{0x000}--\texttt{0x3FF} and data at \texttt{0x400}--\texttt{0x7FF} (the Phase~I baseline
layout) and answers each request after $L$ cycles. We compare three organisations:
\begin{itemize}
  \item \textbf{Modified Harvard} (\sig{mh\_top}): \sig{cpu\_core} with an $N$-line
        instruction cache and an $N$-line data cache, which share \sig{main\_mem} through an
        arbiter.
  \item \textbf{Unified cache}: \sig{cpu\_core\_vn} with one cache of $2N$ lines, the same
        total capacity (\sig{vnc\_top} with \sig{CACHE\_EN}\,=\,1).
  \item \textbf{No cache}: \sig{cpu\_core\_vn} with every access going to \sig{main\_mem}
        (\sig{vnc\_top} with \sig{CACHE\_EN}\,=\,0). With $L=1$ it is the Phase~I baseline,
        cycle for cycle.
\end{itemize}
The Phase~I \sig{harvard\_top}, which has ideal single-cycle memories, remains the reference.
In both Phase~II tops the address decoder sends only data-RAM accesses to a cache.
\sig{out\_port} and invalid addresses bypass the caches and finish in their own cycle.

\begin{figure*}[t]
  \centering
  \resizebox{0.80\textwidth}{!}{% Phase II block diagrams: (a) modified Harvard (mh_top), (b) single bus (vnc_top)
\begin{tikzpicture}[
    font=\footnotesize,
    blk/.style={draw, thick, rounded corners=2pt, minimum height=8mm, align=center, fill=white},
    mem/.style={blk, fill=gray!12},
    cac/.style={blk, fill=blue!8},
    arr/.style={-{Stealth[length=2mm]}, thick},
    lbl/.style={font=\scriptsize\ttfamily, inner sep=1.5pt}]
  % ---------------- (a) modified Harvard ----------------
  \node[blk, minimum width=18mm, minimum height=30mm] (core) at (0,0) {\sig{cpu\_core}\\single\\cycle};
  \node[cac, minimum width=18mm, minimum height=10mm] (ic) at (4.0,0.85) {I-cache\\$N$ lines};
  \node[blk, minimum width=16mm, minimum height=10mm] (dec) at (3.2,-0.85) {decoder\\\sig{rw\_ctrl}};
  \node[cac, minimum width=17mm, minimum height=12mm] (dc) at (5.7,-0.85) {D-cache\\$N$ lines\\+ wbuf};
  \node[blk, minimum width=15mm, minimum height=30mm] (arb) at (8.2,0) {arbiter\\I $>$ D\\no pre-\\emption};
  \node[mem, minimum width=17mm, minimum height=30mm] (mm) at (10.8,0) {\sig{main\_mem}\\code +\\data\\latency $L$};
  \node[blk, minimum width=14mm, minimum height=6mm, font=\scriptsize] (io) at (3.2,-2.3) {\sig{out\_port}};
  \draw[arr] (0.9,1.10) -- node[lbl, above] {i\_addr} (3.1,1.10);
  \draw[arr] (3.1,0.60) -- node[lbl, below] {i\_ready} (0.9,0.60);
  \draw[arr] (0.9,-0.60) -- node[lbl, above] {d\_addr} (2.4,-0.60);
  \draw[arr] (2.4,-1.10) -- node[lbl, below] {d\_ready} (0.9,-1.10);
  \draw[arr] (dec.east) -- node[lbl, above] {RAM} (dc.west);
  \draw[arr] (dec.south) -- node[lbl, right] {IO} (io.north);
  \draw[arr] (ic.east) -- node[lbl, above] {I} (ic.east -| arb.west);
  \draw[arr] (dc.east) -- node[lbl, above] {D} (dc.east -| arb.west);
  \draw[arr, <->] (arb.east) -- node[lbl, above=3pt] {1 port} (mm.west);
  \node[font=\small\bfseries, anchor=west] at (-0.9,1.95) {(a) Modified Harvard (\sig{mh\_top})};
  % ---------------- (b) single bus ----------------
  \begin{scope}[yshift=-4.6cm]
    \node[blk, minimum width=18mm, minimum height=18mm] (vcore) at (0,0) {\sig{cpu\_core\_vn}\\FETCH /\\MEM};
    \node[blk, minimum width=16mm] (vdec) at (3.4,0) {decoder};
    \node[cac, minimum width=24mm, minimum height=12mm] (uc) at (7.0,0) {unified cache\\$2N$ lines + wbuf\\(or uncached)};
    \node[mem, minimum width=17mm, minimum height=18mm] (vmm) at (10.8,0) {\sig{main\_mem}\\code +\\data, $L$};
    \node[blk, minimum width=14mm, minimum height=6mm, font=\scriptsize] (vio) at (3.4,-1.4) {\sig{out\_port}};
    \draw[arr, <->] (vcore.east) -- node[lbl, above] {m\_addr} node[lbl, below] {one bus} (vdec.west);
    \draw[arr, <->] (vdec.east) -- node[lbl, above] {code, RAM} (uc.west);
    \draw[arr] (vdec.south) -- node[lbl, right] {IO} (vio.north);
    \draw[arr, <->] (uc.east) -- node[lbl, above=3pt] {1 port} (vmm.west);
    \node[font=\small\bfseries, anchor=west] at (-0.9,1.35) {(b) Single bus: unified cache or no cache (\sig{vnc\_top})};
  \end{scope}
\end{tikzpicture}
}
  \caption{Phase II. (a) Modified Harvard: split L1 caches over one shared memory, with an
  arbiter giving instruction refills priority. (b) Single bus: the same memory behind a unified
  cache of twice the size, or behind no cache. Both systems use the same cache module and the
  same \sig{main\_mem}.}
  \label{fig:block2}
\end{figure*}

\subsection{Processor extensions}
To run realistic programs, both processors gained XOR, SLT, SLL and SRL, their immediate forms
XORI, SLTI, SLLI and SRLI, the branches BLT and BGE, the jumps JAL and JALR, and LUI, all with
standard RV32I encodings \cite{riscv_spec}. The ALU gained four operations (XOR, signed
set-less-than, and the two logical shifts by \sig{b[4:0]}). BLT and BGE reuse the zero flag after
an SLT, the way BEQ and BNE use it after a SUB. JAL and JALR write $\mathrm{PC}+4$, LUI writes the
U-immediate, and the immediate generator gained the U and J formats. The decoder keeps every
Phase~I output with its old meaning, so the Phase~I subset decodes exactly as before.

Each core also gained a stall interface:
\begin{itemize}
  \item \sig{i\_ready}\,=\,0 (\sig{cpu\_core}): no valid instruction yet. The core decodes a
        NOP, so nothing is written or accessed, and holds its PC.
  \item \sig{d\_ready}\,=\,0 (\sig{cpu\_core}): the load or store of this cycle has not
        finished. The PC and register file are held and the data request stays on the port.
  \item \sig{m\_ready}\,=\,0 (\sig{cpu\_core\_vn}): the bus transfer has not finished. The
        state, PC, latched instruction and register file are held.
\end{itemize}
The Phase~I top levels tie these inputs to 1. With them tied, the enable equations reduce to the
Phase~I ones, and re-running the Phase~I regression gives byte-identical logs: every Phase~I
cycle count is unchanged.

\subsection{Timing design rules}
Every path from the memory back to the processor must end in a register. Otherwise the system
contains a combinational loop, which a static timing tool cannot analyse. We followed two rules:
\begin{enumerate}
  \item \textbf{No forwarding of refill data.} When a miss is served, the cache writes the line
        into its array and answers as a hit in the \emph{next} cycle. If the refill data went
        straight to the core, it would feed the decoder, the ALU, the next data address, the
        D-cache request and the arbiter's address multiplexer, and so reach the memory input
        again.
  \item \textbf{The I-cache request is \sig{\textasciitilde halted}, a register.} If it were
        \sig{\textasciitilde halt}, which also depends on the EBREAK decoded from the fetched
        word, the request would depend on the cache's own output.
\end{enumerate}
Section~\ref{sec:area} shows that the modified Harvard design synthesises with no loop warning,
and that copies which break either rule produce loops.
